\documentclass[
  man,
  floatsintext,
  longtable,
  nolmodern,
  notxfonts,
  notimes,
  draftfirst,
  colorlinks=true,linkcolor=blue,citecolor=blue,urlcolor=blue]{apa7}

\usepackage{amsmath}
\usepackage{amssymb}




\RequirePackage{longtable}
\RequirePackage{threeparttablex}

\makeatletter
\renewcommand{\paragraph}{\@startsection{paragraph}{4}{\parindent}%
	{0\baselineskip \@plus 0.2ex \@minus 0.2ex}%
	{-.5em}%
	{\normalfont\normalsize\bfseries\typesectitle}}

\renewcommand{\subparagraph}[1]{\@startsection{subparagraph}{5}{0.5em}%
	{0\baselineskip \@plus 0.2ex \@minus 0.2ex}%
	{-\z@\relax}%
	{\normalfont\normalsize\bfseries\itshape\hspace{\parindent}{#1}\textit{\addperi}}{\relax}}
\makeatother




\usepackage{longtable, booktabs, multirow, multicol, colortbl, hhline, caption, array, float, xpatch}
\setcounter{topnumber}{2}
\setcounter{bottomnumber}{2}
\setcounter{totalnumber}{4}
\renewcommand{\topfraction}{0.85}
\renewcommand{\bottomfraction}{0.85}
\renewcommand{\textfraction}{0.15}
\renewcommand{\floatpagefraction}{0.7}

\usepackage{tcolorbox}
\tcbuselibrary{listings,theorems, breakable, skins}
\usepackage{fontawesome5}

\definecolor{quarto-callout-color}{HTML}{909090}
\definecolor{quarto-callout-note-color}{HTML}{0758E5}
\definecolor{quarto-callout-important-color}{HTML}{CC1914}
\definecolor{quarto-callout-warning-color}{HTML}{EB9113}
\definecolor{quarto-callout-tip-color}{HTML}{00A047}
\definecolor{quarto-callout-caution-color}{HTML}{FC5300}
\definecolor{quarto-callout-color-frame}{HTML}{ACACAC}
\definecolor{quarto-callout-note-color-frame}{HTML}{4582EC}
\definecolor{quarto-callout-important-color-frame}{HTML}{D9534F}
\definecolor{quarto-callout-warning-color-frame}{HTML}{F0AD4E}
\definecolor{quarto-callout-tip-color-frame}{HTML}{02B875}
\definecolor{quarto-callout-caution-color-frame}{HTML}{FD7E14}

%\newlength\Oldarrayrulewidth
%\newlength\Oldtabcolsep


\usepackage{hyperref}




\providecommand{\tightlist}{%
  \setlength{\itemsep}{0pt}\setlength{\parskip}{0pt}}
\usepackage{longtable,booktabs,array}
\usepackage{calc} % for calculating minipage widths
% Correct order of tables after \paragraph or \subparagraph
\usepackage{etoolbox}
\makeatletter
\patchcmd\longtable{\par}{\if@noskipsec\mbox{}\fi\par}{}{}
\makeatother
% Allow footnotes in longtable head/foot
\IfFileExists{footnotehyper.sty}{\usepackage{footnotehyper}}{\usepackage{footnote}}
\makesavenoteenv{longtable}

\usepackage{graphicx}
\makeatletter
\newsavebox\pandoc@box
\newcommand*\pandocbounded[1]{% scales image to fit in text height/width
  \sbox\pandoc@box{#1}%
  \Gscale@div\@tempa{\textheight}{\dimexpr\ht\pandoc@box+\dp\pandoc@box\relax}%
  \Gscale@div\@tempb{\linewidth}{\wd\pandoc@box}%
  \ifdim\@tempb\p@<\@tempa\p@\let\@tempa\@tempb\fi% select the smaller of both
  \ifdim\@tempa\p@<\p@\scalebox{\@tempa}{\usebox\pandoc@box}%
  \else\usebox{\pandoc@box}%
  \fi%
}
% Set default figure placement to htbp
\def\fps@figure{htbp}
\makeatother


% definitions for citeproc citations
\NewDocumentCommand\citeproctext{}{}
\NewDocumentCommand\citeproc{mm}{%
  \begingroup\def\citeproctext{#2}\cite{#1}\endgroup}
\makeatletter
 % allow citations to break across lines
 \let\@cite@ofmt\@firstofone
 % avoid brackets around text for \cite:
 \def\@biblabel#1{}
 \def\@cite#1#2{{#1\if@tempswa , #2\fi}}
\makeatother
\newlength{\cslhangindent}
\setlength{\cslhangindent}{1.5em}
\newlength{\csllabelwidth}
\setlength{\csllabelwidth}{3em}
\newenvironment{CSLReferences}[2] % #1 hanging-indent, #2 entry-spacing
 {\begin{list}{}{%
  \setlength{\itemindent}{0pt}
  \setlength{\leftmargin}{0pt}
  \setlength{\parsep}{0pt}
  % turn on hanging indent if param 1 is 1
  \ifodd #1
   \setlength{\leftmargin}{\cslhangindent}
   \setlength{\itemindent}{-1\cslhangindent}
  \fi
  % set entry spacing
  \setlength{\itemsep}{#2\baselineskip}}}
 {\end{list}}
\usepackage{calc}
\newcommand{\CSLBlock}[1]{\hfill\break\parbox[t]{\linewidth}{\strut\ignorespaces#1\strut}}
\newcommand{\CSLLeftMargin}[1]{\parbox[t]{\csllabelwidth}{\strut#1\strut}}
\newcommand{\CSLRightInline}[1]{\parbox[t]{\linewidth - \csllabelwidth}{\strut#1\strut}}
\newcommand{\CSLIndent}[1]{\hspace{\cslhangindent}#1}





\usepackage{newtx}

\defaultfontfeatures{Scale=MatchLowercase}
\defaultfontfeatures[\rmfamily]{Ligatures=TeX,Scale=1}





\title{Intrusive and Voluntary Memory from a Single System: A
Retrieved-Context Account of Selective Interference}


\shorttitle{Intrusive and Voluntary Memory from a Single System}


\usepackage{etoolbox}






\author{Jordan B Gunn}



\affiliation{
{Department of Psychology, University of Cambridge}}




\leftheader{Gunn}



\abstract{Post-encoding visuospatial tasks such as playing Tetris after
a trauma reminder can reduce subsequent intrusive memories while leaving
voluntary memory (e.g., recognition and intentional recall) relatively
intact. This selective interference effect has been taken as evidence
for separate ``intrusive'' versus ``voluntary'' memory traces or for
modality-specific disruption of consolidation or reconsolidation. Here
we develop a single-system account within retrieved-context theory
(CMR3) that derives selective interference from generic episodic binding
and retrieval dynamics. In the model, reminders reinstate
trauma-associated context, and experiences during the interference task
are encoded in that reinstated context, creating many strong competitor
episodes. Later, as ongoing context drifts through the trauma region,
uncontrolled context-to-item retrieval is increasingly captured by these
competitors, reducing the probability that any single trauma memory is
sampled as an intrusion. In contrast, probe-driven item-to-context
retrieval and controlled context-to-item recall can remain largely
intact, explaining spared recognition and (under sufficient control)
intentional recall. Simulations of trauma-film paradigms reproduce the
intrusion--recognition dissociation and its heterogeneity across
designs, and yield testable predictions about timing, cueing, emotional
arousal, and retrieval-control manipulations. This framework reframes
visuospatial tasks as one instantiation of a broader competitor-encoding
intervention strategy and provides principled boundary conditions for
translating lab effects to clinical settings. }

\keywords{episodic memory, intrusive memories, selective
interference, trauma-film paradigm, retrieved-context
theory, computational modeling}

\authornote{ 

\par{  This is an in-progress working draft. Please do not cite, quote,
or circulate without the authors' written permission.     }
}

\makeatletter
\let\endoldlt\endlongtable
\def\endlongtable{
\hline
\endoldlt
}
\makeatother

\urlstyle{same}



\makeatletter
\@ifpackageloaded{caption}{}{\usepackage{caption}}
\AtBeginDocument{%
\ifdefined\contentsname
  \renewcommand*\contentsname{Table of contents}
\else
  \newcommand\contentsname{Table of contents}
\fi
\ifdefined\listfigurename
  \renewcommand*\listfigurename{List of Figures}
\else
  \newcommand\listfigurename{List of Figures}
\fi
\ifdefined\listtablename
  \renewcommand*\listtablename{List of Tables}
\else
  \newcommand\listtablename{List of Tables}
\fi
\ifdefined\figurename
  \renewcommand*\figurename{Figure}
\else
  \newcommand\figurename{Figure}
\fi
\ifdefined\tablename
  \renewcommand*\tablename{Table}
\else
  \newcommand\tablename{Table}
\fi
}
\@ifpackageloaded{float}{}{\usepackage{float}}
\floatstyle{ruled}
\@ifundefined{c@chapter}{\newfloat{codelisting}{h}{lop}}{\newfloat{codelisting}{h}{lop}[chapter]}
\floatname{codelisting}{Listing}
\newcommand*\listoflistings{\listof{codelisting}{List of Listings}}
\makeatother
\makeatletter
\makeatother
\makeatletter
\@ifpackageloaded{caption}{}{\usepackage{caption}}
\@ifpackageloaded{subcaption}{}{\usepackage{subcaption}}
\makeatother

% From https://tex.stackexchange.com/a/645996/211326
%%% apa7 doesn't want to add appendix section titles in the toc
%%% let's make it do it
\makeatletter
\xpatchcmd{\appendix}
  {\par}
  {\addcontentsline{toc}{section}{\@currentlabelname}\par}
  {}{}
\makeatother

%% Disable longtable counter
%% https://tex.stackexchange.com/a/248395/211326

\usepackage{etoolbox}

\makeatletter
\patchcmd{\LT@caption}
  {\bgroup}
  {\bgroup\global\LTpatch@captiontrue}
  {}{}
\patchcmd{\longtable}
  {\par}
  {\par\global\LTpatch@captionfalse}
  {}{}
\apptocmd{\endlongtable}
  {\ifLTpatch@caption\else\addtocounter{table}{-1}\fi}
  {}{}
\newif\ifLTpatch@caption
\makeatother

\begin{document}

\maketitle


\setcounter{secnumdepth}{-\maxdimen} % remove section numbering

\setlength\LTleft{0pt}


\section{Introduction}\label{introduction}

Traumatic experience foregrounds a trade-off between durability and
selectivity in episodic remembering. Memory should retain highly
consequential episodes so critical details stay accessible long into the
future. At the same time, adaptive retrieval is selective: what comes to
mind should normally depend on current cues, goals, and safety.
Intrusive memories after trauma suggest that durable encoding can
outpace regulatory control. The event returns repeatedly and
indiscriminately, despite the demands of the present
(\citeproc{ref-brewin2010intrusive}{Brewin et al., 2010};
\citeproc{ref-ehlers2000cognitive}{Ehlers \& Clark, 2000};
\citeproc{ref-iyadurai2019intrusive}{Iyadurai et al., 2019};
\citeproc{ref-krans2009intrusive}{Krans et al., 2009}).

The present work examines this trade-off in the trauma-film paradigm
(\citeproc{ref-holmes2008inducing}{Holmes \& Bourne, 2008};
\citeproc{ref-james2016trauma}{James et al., 2016}). Participants watch
a distressing film depicting scenes of injury or threat
(\citeproc{ref-holmes2008inducing}{Holmes \& Bourne, 2008};
\citeproc{ref-james2016trauma}{James et al., 2016}). Intrusive memories
commonly emerge over the following days in diary measures
(\citeproc{ref-holmes2004trauma}{Holmes et al., 2004};
\citeproc{ref-holmes2008inducing}{Holmes \& Bourne, 2008};
\citeproc{ref-james2016trauma}{James et al., 2016}). When asked to
recall or recognize the same material deliberately, voluntary memory can
be relatively preserved, and its relationship to intrusion frequency can
be weak, such that the two measures dissociate
(\citeproc{ref-holmes2004trauma}{Holmes et al., 2004};
\citeproc{ref-james2016trauma}{James et al., 2016}). The same encoded
episode can therefore support both uncontrolled and controlled access,
and these modes can come apart empirically
(\citeproc{ref-holmes2004trauma}{Holmes et al., 2004};
\citeproc{ref-iyadurai2019intrusive}{Iyadurai et al., 2019}).

A sharper dissociation emerges when a brief visuospatial task follows
the film or a later film reminder (\citeproc{ref-holmes2009can}{Holmes
et al., 2009}; \citeproc{ref-james2015computer}{James et al., 2015};
\citeproc{ref-kessler2020visuospatial}{Kessler et al., 2020}). Across
studies, participants who play Tetris or perform a similarly
visually-engaging task report fewer intrusive memories over the
following days (\citeproc{ref-deeprose2012imagery}{Deeprose et al.,
2012}; \citeproc{ref-holmes2009can}{Holmes et al., 2009},
\citeproc{ref-holmes2010key}{2010}). Voluntary memory for the same
material generally shows little measurable impact even when retrieval
conditions are matched
(\citeproc{ref-asselbergs2023systematic}{Asselbergs et al., 2023};
\citeproc{ref-lau2019intrusive}{Lau-Zhu et al., 2019}). Such
interventions are effective even after prompting reinstatement of
real-life traumatic memories, reducing intrusions while preserving
deliberate recall for life experiences such as motor vehicle accidents,
traumatic childbirth following emergency caesarean section, and
war-related trauma in refugees (\citeproc{ref-horsch2017reducing}{Horsch
et al., 2017}; \citeproc{ref-iyadurai2018preventing}{Iyadurai et al.,
2018}; \citeproc{ref-kanstrup2021single}{Kanstrup et al., 2021}). This
selective interference effect deepens the dissociation. A post-encoding
or post-reminder manipulation appears to alter one mode of retrieval
more than the other, even when both draw on the same recent experience.

A prominent interpretation attributes this selectivity to separable
memory representations (\citeproc{ref-brewin1996dual}{Brewin et al.,
1996}; \citeproc{ref-brewin2014episodic}{Brewin, 2014}). Under the
dual-representation account, traumatic experience produces distinct
sensory and contextual memory traces --- the first supporting
involuntary re-experiencing, the second supporting deliberate recall
(\citeproc{ref-brewin2010intrusive}{Brewin et al., 2010};
\citeproc{ref-brewin2014episodic}{Brewin, 2014}). A visuospatial
interference task competes for the perceptual processing resources
needed to consolidate the sensory trace, reducing later intrusions while
leaving the contextual trace and the voluntary memory it supports
relatively intact (\citeproc{ref-holmes2009can}{Holmes et al., 2009};
\citeproc{ref-james2015computer}{James et al., 2015}). When the
intervention follows a reminder at longer delays, the account invokes
reconsolidation: the reminder reactivates the sensory trace, and the
interfering task disrupts its restabilization during a time-limited
window (\citeproc{ref-james2015computer}{James et al., 2015};
\citeproc{ref-kindt2009beyond}{Kindt et al., 2009}). The account thus
ties selective interference to separate memory representations,
modality-specific disruption of one trace, and --- for delayed effects
--- a discrete reconsolidation window.

Here we develop an alternative within a unitary retrieved-context
framework and demonstrate it through simulations of a formal model. The
framework holds that items are bound to a continuously evolving temporal
context that later serves as the primary retrieval cue
(\citeproc{ref-howard2002distributed}{Howard \& Kahana, 2002};
\citeproc{ref-polyn2009context}{Polyn et al., 2009}). In this account,
the original event and a later interference task deposit traces in this
shared context, and the degree of overlap between their contexts governs
retrieval competition. Retrieval control mechanisms available during
deliberate recall, reinstating the event's context and sharpening
competition among candidates, can preserve the target's advantage over
competitors. The dissociation between involuntary and voluntary recall
emerges from these interacting constraints without requiring separate
traces, reconsolidation windows, or modality-specific disruption as
modeling prerequisites. By grounding selective interference in general
principles of temporal context and retrieval competition, the account
connects a clinical phenomenon to the broader episodic memory literature
and generates predictions that diverge from the dual-representation
account and can sharpen intervention designs.

\section{Overview of the Model}\label{overview-of-the-model}

The account developed here builds on a contextual binding theory of
episodic memory (\citeproc{ref-howard2002distributed}{Howard \& Kahana,
2002}; \citeproc{ref-yonelinas2019contextual}{Yonelinas et al., 2019}).
In this framework, each experienced event is encoded by binding it to
the current state of a continuously evolving context representation.
Context encompasses temporal, spatial, and cognitive features, some of
which change rapidly from one event to the next while others shift
gradually with ongoing experience
(\citeproc{ref-yonelinas2019contextual}{Yonelinas et al., 2019}).
Forgetting, on this view, arises primarily from contextual interference
rather than time-dependent decay or consolidation. Events encoded in
similar context states compete during retrieval, and the more
competitors share a target's contextual neighborhood, the less likely
the target is to be retrieved (\citeproc{ref-mensink1989model}{Mensink
\& Raaijmakers, 1989}; \citeproc{ref-yonelinas2019contextual}{Yonelinas
et al., 2019}).

Retrieved-context theory formalizes this binding framework with two
associative pathways (\citeproc{ref-howard2002distributed}{Howard \&
Kahana, 2002}; \citeproc{ref-polyn2009context}{Polyn et al., 2009}).
Context-to-item associations allow the current context state to cue all
events encoded in similar states, producing competitive retrieval in
which events vie for recall based on the strength of their contextual
match. Item-to-context associations allow a recalled or re-encountered
event to reinstate the context in which it was originally experienced, a
mechanism absent from earlier context theories. This reinstatement
mechanism creates a bidirectional loop: retrieving an event drives
context back toward the state present at encoding, which in turn cues
neighboring events, producing the temporal contiguity effect
(\citeproc{ref-kahana1996associative}{Kahana, 1996};
\citeproc{ref-sederberg2008context}{Sederberg et al., 2008}).

The critical property for the present account is that context-to-item
retrieval is competitive: any event encoded in the current contextual
neighborhood is a candidate, and adding new events to that neighborhood
increases retrieval competition. In the trauma-film paradigm, intrusive
memories arise through this pathway as ongoing context drifts near
states associated with the film, cueing film events without intention.
When a participant performs a visuospatial task such as Tetris
immediately after viewing a distressing film, the events the task
produces are encoded in temporal context that overlaps substantially
with the film's contextual neighborhood. When context later drifts near
the film region, interference events compete alongside film events,
reducing the probability that any single film event is retrieved as an
intrusion. The strength of this interference depends on how many
competitors are encoded and how much their context overlaps with the
film's, not on whether the task shares the film's perceptual modality.

The item-to-context pathway extends this account to delayed
interventions. In the trauma-film paradigm, presenting a film clip as a
reminder retrieves its associated film context via item-to-context
retrieval, pulling the current context back toward the state present
during the original viewing (\citeproc{ref-james2015computer}{James et
al., 2015}; \citeproc{ref-kessler2020visuospatial}{Kessler et al.,
2020}). In clinical interventions at longer delays, directed
recollection of a real-life traumatic experience serves the same
function, reinstating the original trauma context through the same
associative pathway (\citeproc{ref-horsch2017reducing}{Horsch et al.,
2017}; \citeproc{ref-iyadurai2018preventing}{Iyadurai et al., 2018}). A
visuospatial task performed in this reinstated state encodes competitors
that overlap with the trauma context, enabling interference through the
same contextual competition as the immediate case. Interference
effectiveness depends on two continuous variables: how strongly the
reminder reinstates film context, and how much subsequent experience has
caused the reinstated context to drift before the interference task
begins. Both are governed by context overlap rather than a discrete
reconsolidation window that opens and closes.

Item-to-context associations also create a retrieval pathway that is
structurally immune to contextual competition. When a film event is
presented as a recognition probe, it retrieves its own associated
context through its item-to-context associations. Competitor events
encoded after the film modify only their own associations; they do not
alter the film event's stored context. Recognition therefore bypasses
the competitive dynamics of context-to-item retrieval entirely. This
architectural immunity contrasts with the graded protection that
retrieval control provides to deliberate recall, a distinction developed
below.

Deliberate recall of the film presents a subtler case. Like intrusions,
it uses context-to-item retrieval and is therefore exposed to competitor
interference. Two control mechanisms can partially offset this
vulnerability. Reinstating the film's starting context as the initial
retrieval cue biases recall toward the film region and away from the
interference neighborhood; sharpening competition among candidates
amplifies strongly associated film events relative to weakly matched
competitors. Together, these mechanisms produce graded protection. This
protection falls short of recognition's architectural immunity but
meaningfully reduces vulnerability, with the degree of sparing depending
on how much control the participant exerts.

These mechanisms reproduce several core findings from the episodic
memory literature. In free recall, the context-reinstatement loop
accounts for temporal contiguity (the tendency to recall temporally
neighboring events in sequence) and for the persistence of this pattern
across long filled delays (\citeproc{ref-howard2002distributed}{Howard
\& Kahana, 2002}; \citeproc{ref-kahana1996associative}{Kahana, 1996};
\citeproc{ref-sederberg2008context}{Sederberg et al., 2008}). Memories
from recent but irrelevant experiences also intrude during recall of a
target episode, with a frequency that decays gradually over intervening
time (\citeproc{ref-howard2008persistence}{Howard et al., 2008};
\citeproc{ref-lohnas2015expanding}{Lohnas et al., 2015};
\citeproc{ref-zaromb2006temporal}{Zaromb et al., 2006}). Extensions of
the model (eCMR, CMR3) add emotional arousal and valence as dimensions
of context, so that high-arousal events are bound not only to temporal
context but also to an emotional context reflecting the affective state
at encoding (\citeproc{ref-talmi2019retrieved}{Talmi et al., 2019}).
This extension is directly relevant to the present work, as traumatic
film material is encoded with strong arousal, and the model's treatment
of arousal as context captures the enhanced encoding and distinctive
retrieval dynamics that characterize emotional material.

This account replaces each of the three commitments that the
dual-representation account requires. The dissociation between
involuntary and voluntary retrieval emerges from two retrieval pathways
operating on a single set of memory associations, not from separate
sensory and contextual traces. Interference depends on competitor
encoding strength and context overlap, not on modality-specific
disruption of a sensory trace. Delayed effects follow from context
reinstatement and its maintenance across intervening experience, rather
than from a time-limited reconsolidation window.

These mechanisms jointly predict a gradient of vulnerability to
interference. Recognition should be fully spared because it uses a
pathway that is architecturally immune to competitor encoding.
Deliberate recall should be partially spared because control mechanisms
can offset competition, with the degree of protection varying with the
control resources available. Unguided recall, the pathway that produces
intrusions, should be most vulnerable because it lacks both
architectural immunity and active control. This graded immunity
prediction falls out of the model's architecture without additional
assumptions and cannot be produced by an account in which the
dissociation arises from separate memory stores, one disrupted and one
intact. The simulations that follow test this prediction across a series
of trauma-film paradigm designs.

\section{Model and Paradigm}\label{model-and-paradigm}

\subsection{Representations and Memory
Architecture}\label{representations-and-memory-architecture}

The model represents each study event as an orthonormal feature vector
\(f_i\). In the trauma-film paradigm, events span four experimental
phases presented in fixed order: film clips, break activities,
interference items, and filler items, for a total of \(N\) items. Each
item carries an emotional flag \(e_i \in \{0, 1\}\) indicating whether
it was encoded under heightened arousal.

Two context representations evolve across the experiment. Temporal
context \(c^T\) is a unit-length vector of dimension \(N + 1\),
initialized with all activation on a dedicated start-of-list unit. As
each item is encoded, \(c^T\) drifts toward the item's retrieved
context. The result is a recency-weighted trace of recent experience
that serves as the primary retrieval cue. Emotional context \(c^E\) is a
separate unit-length vector with two dimensions --- a neutral pole and
an arousal pole --- initialized at \(c^E_0 = [1, 0]\). When emotional
items are encoded, \(c^E\) drifts toward the arousal pole, creating
shared emotional context among co-aroused items.

Four associative memory matrices link items to contexts. Item-to-context
memory \(M^{FC}\) stores associations from item features to temporal
context and supports context reinstatement when an item is retrieved or
re-encountered. Context-to-item memory \(M^{CF}\) stores the reverse
associations and supports competitive retrieval: the current temporal
context cues all items encoded in similar states. Two parallel matrices
serve the emotional pathway: \(M^{FC}_E\) maps item features to
emotional context, and \(M^{CF}_E\) maps emotional context to item
features.

Table~\ref{tbl-parameters} lists all model structures and parameters;
these representations and their initializations underpin every
simulation in this paper.

\begin{table}

{\caption{{Model structures and parameters. Base parameters
(\(\beta_{enc}\) through \(\theta_r\), excluding \(\beta_{emot}\) and
\(\omega_E\)) are fitted per-subject to the free-recall benchmark of
Healey and Kahana
(\citeproc{ref-healey2014memory}{2014}).}{\label{tbl-parameters}}}
\vspace{-20pt}}

\begin{longtable}[]{@{}
  >{\raggedright\arraybackslash}p{(\linewidth - 4\tabcolsep) * \real{0.3333}}
  >{\raggedright\arraybackslash}p{(\linewidth - 4\tabcolsep) * \real{0.3333}}
  >{\raggedright\arraybackslash}p{(\linewidth - 4\tabcolsep) * \real{0.3333}}@{}}
\toprule\noalign{}
\begin{minipage}[b]{\linewidth}\raggedright
Symbol
\end{minipage} & \begin{minipage}[b]{\linewidth}\raggedright
Name
\end{minipage} & \begin{minipage}[b]{\linewidth}\raggedright
Description
\end{minipage} \\
\midrule\noalign{}
\endhead
\bottomrule\noalign{}
\endlastfoot
\(c^T\) & temporal context & Recency-weighted summary of recent items \\
\(c^E\) & emotional context & Arousal state at encoding (2-D) \\
\(f_i\) & item features & Orthonormal representation of item \(i\) \\
\(e_i\) & emotional flag & 1 for emotional items, 0 for neutral \\
\(M^{FC}\), \(M^{CF}\) & temporal memories & Item \(\leftrightarrow\)
temporal-context associations \\
\(M^{FC}_E\), \(M^{CF}_E\) & emotional memories & Item
\(\leftrightarrow\) emotional-context associations \\
\(\beta_{enc}\) & encoding drift rate & Temporal context drift during
study \\
\(\beta_{start}\) & start drift rate & Start-of-list reinstatement at
recall onset \\
\(\beta_{rec}\) & recall drift rate & Context drift during recall \\
\(\beta_{emot}\) & emotional drift rate & Emotional context drift during
encoding \\
\(\gamma\) & learning rate & Feature-to-context learning rate \\
\(\alpha\) & shared support & Pre-experimental shared support in
\(M^{CF}\) \\
\(\delta\) & item support & Pre-experimental self-support in \(M^{CF}\)
and \(M^{CF}_E\) \\
\(\phi_s\) & primacy scale & Initial boost to context-to-item
learning \\
\(\phi_d\) & primacy decay & Decay rate of the primacy boost \\
\(\tau\) & choice sensitivity & Exponent for Luce-style retrieval
competition \\
\(\omega_E\) & emotional weight & Strength of emotional pathway in
activation \\
\(\theta_s\) & stop scale & Baseline termination probability \\
\(\theta_r\) & stop growth & Termination growth with recall position \\
\end{longtable}

\end{table}

Pre-experimental initialization establishes baseline associations before
the experiment. For the temporal pathway:

\[
M^{FC}_{pre(ij)} =
\begin{cases}
1 - \gamma & \text{if } j = i \\
0 & \text{otherwise}
\end{cases}
\]

\[
M^{CF}_{pre(ij)} =
\begin{cases}
\delta & \text{if } i = j \\
\alpha & \text{otherwise}
\end{cases}
\]

Here, \(\gamma\) scales the ratio of experimental to pre-experimental
item-to-context associations, \(\delta\) sets each item's
pre-experimental self-support, and \(\alpha\) sets shared support across
all items.

For the emotional pathway, pre-experimental associations link emotional
items to the arousal pole of emotional context:

\[
M^{FC}_{E,pre(ik)} =
\begin{cases}
(1 - \gamma) \cdot e_i & \text{if } k = 1 \text{ (arousal)} \\
0 & \text{if } k = 0 \text{ (neutral)}
\end{cases}
\]

\[
M^{CF}_{E,pre(kj)} =
\begin{cases}
\delta \cdot e_j & \text{if } k = 1 \text{ (arousal)} \\
0 & \text{if } k = 0 \text{ (neutral)}
\end{cases}
\]

Here, \(k\) indexes the two emotional context units. Emotional items
(\(e_i = 1\)) establish pre-experimental links through the arousal pole,
creating a channel for emotional retrieval competition distinct from
temporal proximity. Neutral items have no emotional associations.

\subsection{Encoding: Film, Interference, and Filler
Phases}\label{encoding-film-interference-and-filler-phases}

All four experimental phases --- film, break, interference, and filler
--- use the same encoding machinery, differing only in phase-specific
rates. On each encoding step, the model retrieves context input, updates
both context streams, and strengthens item--context associations in all
four memory matrices.

Temporal context input for item \(i\) is retrieved by probing
\(M^{FC}\):

\[
c^{IN}_i = M^{FC} f_i
\]

The model normalizes the input to unit length. Temporal context then
drifts toward the normalized input:

\[
c^T_i = \rho_i \, c^T_{i-1} + \beta_{enc} \, c^{IN}_i
\]

\[
\rho_i = \sqrt{1 + \beta_{enc}^2 \left[\left(c^T_{i-1} \cdot c^{IN}_i\right)^2 - 1\right]} - \beta_{enc} \left(c^T_{i-1} \cdot c^{IN}_i\right)
\]

Here, \(\beta_{enc}\) controls how rapidly temporal context drifts
toward newly encoded items, and \(\rho_i\) preserves unit length.

Feature-to-context associations update with a Hebbian outer product:

\[
\Delta M^{FC}_{ij} = \gamma \, f_i \, c^T_j
\]

Here, \(\gamma\) is the feature-to-context learning rate. A primacy
function that decays with serial position weights context-to-item
associations:

\[
\phi_i = \phi_s \, e^{-\phi_d \cdot i} + 1
\]

\[
\Delta M^{CF}_{ij} = \phi_i \, c^T_j \, f_i
\]

Here, \(\phi_s\) sets the initial primacy boost and \(\phi_d\) controls
its decay. Items encoded early in a phase receive stronger
context-to-item associations, producing a primacy gradient in later
retrieval.

Emotional context updates in parallel. The model retrieves emotional
context input as \(c^{IN}_{E,i} = M^{FC}_E f_i\), normalizes it, and
integrates:

\[
c^E_i = \rho_{E,i} \, c^E_{i-1} + \beta_{emot} \, c^{IN}_{E,i}
\]

For emotional items (\(e_i = 1\)), \(M^{FC}_E\) returns input weighted
toward the arousal pole, causing \(c^E\) to drift toward shared arousal.
For neutral items (\(e_i = 0\)), \(M^{FC}_E\) returns near-zero input,
leaving \(c^E\) essentially unchanged.

Emotional memory learning mirrors the temporal pathway, but the
emotional flag gates it:

\[
\Delta M^{FC}_{E,ij} = \gamma \cdot e_i \, f_i \, c^E_j
\]

\[
\Delta M^{CF}_{E,ij} = \phi_i \cdot e_i \, c^E_j \, f_i
\]

Only emotional items form new emotional associations during encoding.
Neutral events pass through the emotional pathway without modification,
confining emotional retrieval competition to material encoded under
heightened arousal.

Each phase applies multiplicative scale factors to the base encoding
rates. The film phase uses base rates directly. Subsequent phases
modulate two quantities: drift rate, which controls context proximity to
preceding items, and \(M^{CF}\) learning rate, which controls competitor
encoding strength.

\begin{table}

{\caption{{Phase-specific encoding rates. Drift scales multiply
\(\beta_{enc}\) and are clipped to
\([0, 1]\).}{\label{tbl-phase-scales}}}
\vspace{-20pt}}

\begin{longtable}[]{@{}
  >{\raggedright\arraybackslash}p{(\linewidth - 6\tabcolsep) * \real{0.2500}}
  >{\raggedright\arraybackslash}p{(\linewidth - 6\tabcolsep) * \real{0.2500}}
  >{\raggedright\arraybackslash}p{(\linewidth - 6\tabcolsep) * \real{0.2500}}
  >{\raggedright\arraybackslash}p{(\linewidth - 6\tabcolsep) * \real{0.2500}}@{}}
\toprule\noalign{}
\begin{minipage}[b]{\linewidth}\raggedright
Phase
\end{minipage} & \begin{minipage}[b]{\linewidth}\raggedright
Drift rate
\end{minipage} & \begin{minipage}[b]{\linewidth}\raggedright
\(M^{CF}\) learning scale
\end{minipage} & \begin{minipage}[b]{\linewidth}\raggedright
Paradigm role
\end{minipage} \\
\midrule\noalign{}
\endhead
\bottomrule\noalign{}
\endlastfoot
Film & \(\beta_{enc}\) & \(1\) & Target encoding \\
Break & \(s^{\beta}_\text{break} \cdot \beta_{enc}\) &
\(s^{\phi}_\text{break}\) & Retention interval \\
Interference & \(s^{\beta}_\text{int} \cdot \beta_{enc}\) &
\(s^{\phi}_\text{int}\) & Competitor encoding \\
Filler & \(s^{\beta}_\text{fill} \cdot \beta_{enc}\) &
\(s^{\phi}_\text{fill}\) & Pre-recall buffer \\
\end{longtable}

\end{table}

The interference scales \(s^{\beta}_\text{int}\) and
\(s^{\phi}_\text{int}\) are central to the simulations. The drift scale
controls how much temporal context drifts during interference encoding:
higher values produce interference items whose temporal context overlaps
more with the film neighborhood, strengthening competition. The MCF
scale controls competitor encoding strength: stronger competitors
capture more retrieval probability from film targets. The mechanism is
straightforward: because \(M^{CF}\) rows encode the temporal context
present at study, interference items encoded in context that overlaps
the film neighborhood activate whenever the retrieval cue revisits that
region. These two scales jointly define the intervention's competitive
impact and are the primary parameters varied in the \emph{Selective
Interference Effect} and \emph{What Makes Interference Stronger}
simulations.

\subsection{Reminders and Context
Reinstatement}\label{reminders-and-context-reinstatement}

In the trauma-film paradigm, a film clip presented as a reminder between
the break and interference phases serves to reinstate the film's
encoding context. The model implements this in two steps. First, both
context representations drift toward their respective start-of-list
states:

\[
c^T_\text{rem} = \rho \, c^T + \beta_\text{rem,start} \, c^T_0
\]

\[
c^E_\text{rem} = \rho_E \, c^E + \beta_\text{rem,start} \, c^E_0
\]

Here,
\(\beta_\text{rem,start} = s^{\beta}_\text{rem,start} \cdot \beta_{start}\)
sets how strongly context shifts toward the pre-experimental state
before the reminder arrives. This partial reinstatement moves context
back toward the general region of the film's temporal neighborhood.

Second, each reminder item reinstates its specific temporal and
emotional context through the same item-to-context retrieval used during
recall:

\[
c^{IN}_\text{cue} = M^{FC} f_\text{cue}
\]

\[
c^T \leftarrow \rho \, c^T + \beta_\text{rem} \, c^{IN}_\text{cue}
\]

\[
c^{IN}_{E,\text{cue}} = M^{FC}_E f_\text{cue}
\]

\[
c^E \leftarrow \rho_E \, c^E + \beta_\text{rem} \, c^{IN}_{E,\text{cue}}
\]

Here, \(\beta_\text{rem} = s^{\beta}_\text{rem} \cdot \beta_{enc}\)
controls how strongly the reminder pulls context toward the film's
encoding state. Critically, reminders reinstate context without forming
new associations: \(M^{FC}\), \(M^{CF}\), \(M^{FC}_E\), and \(M^{CF}_E\)
remain unchanged. This distinction mirrors the clinical logic of a
reminder cue --- it reactivates the memory without adding new content.

The reminder mechanism has a direct consequence for interference. By
pulling temporal context back toward the film neighborhood, the reminder
ensures that interference items encoded immediately afterward share
temporal context with the film. Stronger reinstatement
(\(s^{\beta}_\text{rem}\) closer to 1) produces more overlap, and hence
more competition, between film targets and interference competitors.
These dynamics are tested in the \emph{Context Reinstatement and Delayed
Interference} simulations, which sweep \(s^{\beta}_\text{rem,start}\)
and \(s^{\beta}_\text{rem}\) to examine how reinstatement strength
modulates interference effectiveness.

\subsection{Retrieval: Three Modes}\label{retrieval-three-modes}

The model distinguishes three retrieval modes that share the same memory
architecture but differ in pathway, control, and vulnerability to
interference.

\textbf{Context-to-item retrieval} underlies both unguided intrusions
and directed free recall. At the start of recall, both context streams
drift toward their initial states:

\[
c^T_\text{start} = \rho \, c^T_N + \beta_{start} \, c^T_0
\]

\[
c^E_\text{start} = \rho_E \, c^E_N + \beta_{start} \, c^E_0
\]

Here, \(\beta_{start}\) controls how much start-of-list context is
reinstated at retrieval onset. The current temporal and emotional
contexts then cue all items through \(M^{CF}\) and \(M^{CF}_E\):

\[
A_T = M^{CF} \, c^T \odot r
\]

\[
A_E = \omega_E \cdot M^{CF}_E \, c^E \odot r
\]

\[
A = A_T + A_E
\]

Here, \(r\) is a binary mask over recallable items (studied and not yet
recalled), and \(\omega_E\) weights the emotional pathway's
contribution. Because all recallable items contribute to the
denominator, encoding additional competitors in contextually similar
states dilutes each film item's share of retrieval probability. The
choice sensitivity \(\tau\) power-scales the combined activation vector
\(A\). The probability of recalling item \(i\) then follows a Luce
choice rule:

\[
P(i \mid \text{recall}) = \frac{A_i^{\tau}}{\sum_k A_k^{\tau}}
\]

At each retrieval step, the model may terminate with probability
\(p_\text{stop}\) or recall an item with probability
\((1 - p_\text{stop}) \cdot P(i \mid \text{recall})\). Termination
probability grows with recall position:

\[
p_\text{stop}(j) = \theta_s \, e^{\theta_r \cdot j}
\]

Here, \(j\) is the number of items already recalled, \(\theta_s\) sets
the baseline termination probability, and \(\theta_r\) controls its
growth. When an item is recalled, both context streams reinstate its
associated context:

\[
c^T \leftarrow \rho \, c^T + \beta_{rec} \, M^{FC} \, f_\text{recalled}
\]

\[
c^E \leftarrow \rho_E \, c^E + \beta_{rec} \, M^{FC}_E \, f_\text{recalled}
\]

This reinstatement drives context toward the recalled item's encoding
neighborhood, biasing subsequent retrievals toward temporally and
emotionally proximate items. This mechanism produces the contiguity and
emotional clustering effects observed in free recall.

Unguided versus directed retrieval differ in their control parameters,
not their computational pathway. Unguided recall --- the pathway that
produces intrusions --- uses moderate start reinstatement and
competition sharpness. Directed recall uses stronger reinstatement
(\(s^{\beta}_\text{start} \cdot \beta_{start}\)) and sharper competition
(\(s^{\tau} \cdot \tau\)). Together, these bias retrieval toward film
targets because film items, encoded earlier in the sequence, have
temporal context closer to the start-of-list state than interference
items do. These control parameters do not change the retrieval
architecture; they adjust how aggressively the participant's current
context cue favors the target episode. The \emph{What Protects Voluntary
Recall} simulations sweep \(s^{\beta}_\text{start}\) and \(s^{\tau}\) to
quantify how much retrieval control can offset competitor interference.

Item-to-context retrieval underlies recognition. When a film event is
presented as a recognition probe, it retrieves its own associated
temporal context through \(M^{FC}\):

\[
P(\text{recognized} \mid f_i) = P(i \mid f_i \text{ cues context})
\]

The critical property is architectural immunity. Competitor items
encoded during the interference phase modify only their own rows in
\(M^{FC}\) and \(M^{CF}\); they do not alter the film item's stored
associations. Recognition therefore bypasses the competitive dynamics of
context-to-item retrieval entirely. No amount of competitor encoding can
degrade a film item's recognition signal, because the signal depends on
the film item's own item-to-context associations, which interference
leaves untouched. This immunity is a structural consequence of the
model's associative architecture, not a parameter-dependent effect that
could be tuned away. The \emph{Recognition Immunity} simulations verify
this prediction by sweeping interference strength and confirming that
recognition performance remains flat.

\subsection{Calibration and Paradigm
Geometry}\label{calibration-and-paradigm-geometry}

Differential evolution fits eleven base parameters per-subject to the
free-recall data of Healey and Kahana
(\citeproc{ref-healey2014memory}{2014}). This dataset provides recall
sequences from 126 participants across multiple sessions of standard
list-learning. These data yield stable per-subject estimates of
encoding, retrieval, and termination parameters. The fitting procedure
maximizes the likelihood of observed recall sequences under the model.
Two additional parameters --- \(\beta_{emot}\) and \(\omega_E\) ---
control the emotional pathway and default to \(\beta_{enc}\) and 0,
respectively, when emotional context is not manipulated.

Scale factors are not fitted. Each scale factor multiplies a base
parameter to implement a phase-specific encoding or retrieval condition
and is either set to a theoretically motivated value or systematically
varied in parameter sweeps. Table~\ref{tbl-scale-factors} lists all
scale factors with their base parameters and simulation connections.

\begin{table}

{\caption{{Scale factors and their simulation connections. Drift scales
are clipped to \([0, 1]\) after
multiplication.}{\label{tbl-scale-factors}}}
\vspace{-20pt}}

\begin{longtable}[]{@{}
  >{\raggedright\arraybackslash}p{(\linewidth - 6\tabcolsep) * \real{0.2500}}
  >{\raggedright\arraybackslash}p{(\linewidth - 6\tabcolsep) * \real{0.2500}}
  >{\raggedright\arraybackslash}p{(\linewidth - 6\tabcolsep) * \real{0.2500}}
  >{\raggedright\arraybackslash}p{(\linewidth - 6\tabcolsep) * \real{0.2500}}@{}}
\toprule\noalign{}
\begin{minipage}[b]{\linewidth}\raggedright
Scale factor
\end{minipage} & \begin{minipage}[b]{\linewidth}\raggedright
Base parameter
\end{minipage} & \begin{minipage}[b]{\linewidth}\raggedright
Role
\end{minipage} & \begin{minipage}[b]{\linewidth}\raggedright
Simulation connection
\end{minipage} \\
\midrule\noalign{}
\endhead
\bottomrule\noalign{}
\endlastfoot
\(s^{\beta}_\text{break}\) & \(\beta_{enc}\) & Break-phase drift rate &
Calibration \\
\(s^{\phi}_\text{break}\) & \(\phi_i\) & Break-phase MCF learning &
Calibration \\
\(s^{\beta}_\text{rem,start}\) & \(\beta_{start}\) & Context
reinstatement before reminder & Context Reinstatement \\
\(s^{\beta}_\text{rem}\) & \(\beta_{enc}\) & Reminder context drift &
Context Reinstatement \\
\(s^{\beta}_\text{int}\) & \(\beta_{enc}\) & Interference-phase drift
rate & Selective Interference, What Makes Interference Stronger \\
\(s^{\phi}_\text{int}\) & \(\phi_i\) & Interference-phase MCF learning &
Selective Interference, What Makes Interference Stronger \\
\(s^{\beta}_\text{fill}\) & \(\beta_{enc}\) & Filler-phase drift rate &
Calibration \\
\(s^{\phi}_\text{fill}\) & \(\phi_i\) & Filler-phase MCF learning &
Calibration \\
\(s^{\beta}_\text{start}\) & \(\beta_{start}\) & Start-of-recall
reinstatement & What Protects Voluntary Recall \\
\(s^{\tau}\) & \(\tau\) & Retrieval competition sharpness & What
Protects Voluntary Recall \\
\(\omega_E\) & --- & Emotional pathway activation weight & What Makes
Interference Stronger \\
\(s^{\beta}_\text{emot}\) & \(\beta_{emot}\) & Emotional context drift
rate & What Makes Interference Stronger \\
\end{longtable}

\end{table}

The paradigm geometry mirrors the trauma-film experiments simulated in
later sections. Each simulated trial presents \(N_\text{film} = 16\)
film items, \(N_\text{break} = 16\) break items, a reminder phase (one
film clip re-presented without new learning), \(N_\text{int} = 16\)
interference items, and \(N_\text{fill} = 16\) filler items, for a total
list length of \(N = 64\). Recall allows up to 48 retrieval attempts per
trial. Simulations average over 100 replications per subject using the
126 per-subject parameter sets from the Healey and Kahana
(\citeproc{ref-healey2014memory}{2014}) fits.

The primary dependent measure is film-item recall under context-to-item
retrieval. Because retrieval proceeds as a sequence of competitive
samples terminated by a growing stop probability, the expected number of
film items recalled depends on how much retrieval probability
competitors capture at each step. Recognition is operationalized as the
probability of retrieving the correct item given its context, using the
model's item-to-context pathway. These dependent measures map directly
onto the intrusion diaries and recognition tests used in trauma-film
studies. They provide a bridge between the model's formal predictions
and the empirical phenomena the simulations target.

\section{Simulations}\label{simulations}

\subsection{The Selective Interference
Effect}\label{the-selective-interference-effect}

\subsection{Context Reinstatement and Delayed
Interference}\label{context-reinstatement-and-delayed-interference}

\subsection{What Protects Voluntary
Recall}\label{what-protects-voluntary-recall}

\subsection{What Makes Interference
Stronger}\label{what-makes-interference-stronger}

\subsection{Recognition Immunity}\label{recognition-immunity}

\subsection{Test-Phase Cues and Paradigm
Heterogeneity}\label{test-phase-cues-and-paradigm-heterogeneity}

\section{General Discussion}\label{general-discussion}

\section{References}\label{references}

\phantomsection\label{refs}
\begin{CSLReferences}{1}{0}
\bibitem[\citeproctext]{ref-asselbergs2023systematic}
Asselbergs, J., Bentum, J. van, Riper, H., Cuijpers, P., Holmes, E., \&
Sijbrandij, M. (2023). A systematic review and meta-analysis of the
effect of cognitive interventions to prevent intrusive memories using
the trauma film paradigm. \emph{Journal of Psychiatric Research},
\emph{159}, 116--129.

\bibitem[\citeproctext]{ref-brewin2014episodic}
Brewin, C. R. (2014). Episodic memory, perceptual memory, and their
interaction: Foundations for a theory of posttraumatic stress disorder.
\emph{Psychological Bulletin}, \emph{140}(1), 69.

\bibitem[\citeproctext]{ref-brewin1996dual}
Brewin, C. R., Dalgleish, T., \& Joseph, S. (1996). A dual
representation theory of posttraumatic stress disorder.
\emph{Psychological Review}, \emph{103}(4), 670.

\bibitem[\citeproctext]{ref-brewin2010intrusive}
Brewin, C. R., Gregory, J. D., Lipton, M., \& Burgess, N. (2010).
Intrusive images in psychological disorders: Characteristics, neural
mechanisms, and treatment implications. \emph{Psychological Review},
\emph{117}(1), 210.

\bibitem[\citeproctext]{ref-deeprose2012imagery}
Deeprose, C., Zhang, S., DeJong, H., Dalgleish, T., \& Holmes, E. A.
(2012). Imagery in the aftermath of viewing a traumatic film: Using
cognitive tasks to modulate the development of involuntary memory.
\emph{Journal of Behavior Therapy and Experimental Psychiatry},
\emph{43}(2), 758--764.

\bibitem[\citeproctext]{ref-ehlers2000cognitive}
Ehlers, A., \& Clark, D. M. (2000). A cognitive model of posttraumatic
stress disorder. \emph{Behaviour Research and Therapy}, \emph{38}(4),
319--345.

\bibitem[\citeproctext]{ref-healey2014memory}
Healey, M. K., \& Kahana, M. J. (2014). Is memory search governed by
universal principles or idiosyncratic strategies? \emph{Journal of
Experimental Psychology: General}, \emph{143}(2), 575.

\bibitem[\citeproctext]{ref-holmes2008inducing}
Holmes, E. A., \& Bourne, C. (2008). Inducing and modulating intrusive
emotional memories: A review of the trauma film paradigm. \emph{Acta
Psychologica}, \emph{127}(3), 553--566.

\bibitem[\citeproctext]{ref-holmes2004trauma}
Holmes, E. A., Brewin, C. R., \& Hennessy, R. G. (2004). Trauma films,
information processing, and intrusive memory development. \emph{Journal
of Experimental Psychology: General}, \emph{133}(1), 3.

\bibitem[\citeproctext]{ref-holmes2009can}
Holmes, E. A., James, E. L., Coode-Bate, T., \& Deeprose, C. (2009). Can
playing the computer game {``tetris''} reduce the build-up of flashbacks
for trauma? A proposal from cognitive science. \emph{PloS One},
\emph{4}(1), e4153.

\bibitem[\citeproctext]{ref-holmes2010key}
Holmes, E. A., James, E. L., Kilford, E. J., \& Deeprose, C. (2010). Key
steps in developing a cognitive vaccine against traumatic flashbacks:
Visuospatial tetris versus verbal pub quiz. \emph{PloS One},
\emph{5}(11), e13706.

\bibitem[\citeproctext]{ref-horsch2017reducing}
Horsch, A., Vial, Y., Favrod, C., Harari, M. M., Blackwell, S. E.,
Watson, P., Iyadurai, L., Bonsall, M. B., \& Holmes, E. A. (2017).
Reducing intrusive traumatic memories after emergency caesarean section:
A proof-of-principle randomized controlled study. \emph{Behaviour
Research and Therapy}, \emph{94}, 36--47.

\bibitem[\citeproctext]{ref-howard2002distributed}
Howard, M. W., \& Kahana, M. J. (2002). A distributed representation of
temporal context. \emph{Journal of Mathematical Psychology},
\emph{46}(3), 269--299.

\bibitem[\citeproctext]{ref-howard2008persistence}
Howard, M. W., Youker, T. E., \& Venkatadass, V. S. (2008). The
persistence of memory: Contiguity effects across hundreds of seconds.
\emph{Psychonomic Bulletin \& Review}, \emph{15}(1), 58--63.

\bibitem[\citeproctext]{ref-iyadurai2018preventing}
Iyadurai, L., Blackwell, S. E., Meiser-Stedman, R., Watson, P. C.,
Bonsall, M. B., Geddes, J. R., Nobre, A. C., \& Holmes, E. A. (2018).
Preventing intrusive memories after trauma via a brief intervention
involving tetris computer game play in the emergency department: A
proof-of-concept randomized controlled trial. \emph{Molecular
Psychiatry}, \emph{23}(3), 674--682.

\bibitem[\citeproctext]{ref-iyadurai2019intrusive}
Iyadurai, L., Visser, R. M., Lau-Zhu, A., Porcheret, K., Horsch, A.,
Holmes, E. A., \& James, E. L. (2019). Intrusive memories of trauma: A
target for research bridging cognitive science and its clinical
application. \emph{Clinical Psychology Review}, \emph{69}, 67--82.

\bibitem[\citeproctext]{ref-james2015computer}
James, E. L., Bonsall, M. B., Hoppitt, L., Tunbridge, E. M., Geddes, J.
R., Milton, A. L., \& Holmes, E. A. (2015). Computer game play reduces
intrusive memories of experimental trauma via reconsolidation-update
mechanisms. \emph{Psychological Science}, \emph{26}(8), 1201--1215.

\bibitem[\citeproctext]{ref-james2016trauma}
James, E. L., Lau-Zhu, A., Clark, I. A., Visser, R. M., Hagenaars, M.
A., \& Holmes, E. A. (2016). The trauma film paradigm as an experimental
psychopathology model of psychological trauma: Intrusive memories and
beyond. \emph{Clinical Psychology Review}, \emph{47}, 106--142.

\bibitem[\citeproctext]{ref-kahana1996associative}
Kahana, M. J. (1996). Associative retrieval processes in free recall.
\emph{Memory \& Cognition}, \emph{24}(1), 103--109.

\bibitem[\citeproctext]{ref-kanstrup2021single}
Kanstrup, M., Kontio, E., Geranmayeh, A., Olofsdotter Lauri, K., Moulds,
M. L., \& Holmes, E. A. (2021). A single case series using visuospatial
task interference to reduce the number of visual intrusive memories of
trauma with refugees. \emph{Clinical Psychology \& Psychotherapy},
\emph{28}(1), 109--123.

\bibitem[\citeproctext]{ref-kessler2020visuospatial}
Kessler, H., Schmidt, A.-C., James, E. L., Blackwell, S. E., Rauchhaupt,
M. von, Harren, K., Kehyayan, A., Clark, I. A., Sauvage, M., Herpertz,
S., et al. (2020). Visuospatial computer game play after memory reminder
delivered three days after a traumatic film reduces the number of
intrusive memories of the experimental trauma. \emph{Journal of Behavior
Therapy and Experimental Psychiatry}, \emph{67}, 101454.

\bibitem[\citeproctext]{ref-kindt2009beyond}
Kindt, M., Soeter, M., \& Vervliet, B. (2009). Beyond extinction:
Erasing human fear responses and preventing the return of fear.
\emph{Nature Neuroscience}, \emph{12}(3), 256--258.

\bibitem[\citeproctext]{ref-krans2009intrusive}
Krans, J., Näring, G., Becker, E. S., \& Holmes, E. A. (2009). Intrusive
trauma memory: A review and functional analysis. \emph{Applied Cognitive
Psychology: The Official Journal of the Society for Applied Research in
Memory and Cognition}, \emph{23}(8), 1076--1088.

\bibitem[\citeproctext]{ref-lau2019intrusive}
Lau-Zhu, A., Henson, R. N., \& Holmes, E. A. (2019). Intrusive memories
and voluntary memory of a trauma film: Differential effects of a
cognitive interference task after encoding. \emph{Journal of
Experimental Psychology: General}, \emph{148}(12), 2154.

\bibitem[\citeproctext]{ref-lohnas2015expanding}
Lohnas, L., Polyn, S. M., \& Kahana, M. J. (2015). Expanding the scope
of memory search: Modeling intralist and interlist effects in free
recall. \emph{Psychological Review}, \emph{122}(2), 337.

\bibitem[\citeproctext]{ref-mensink1989model}
Mensink, G.-J. M., \& Raaijmakers, J. G. (1989). A model for contextual
fluctuation. \emph{Journal of Mathematical Psychology}, \emph{33}(2),
172--186.

\bibitem[\citeproctext]{ref-polyn2009context}
Polyn, S. M., Norman, K. A., \& Kahana, M. J. (2009). A context
maintenance and retrieval model of organizational processes in free
recall. \emph{Psychological Review}, \emph{116}(1), 129.

\bibitem[\citeproctext]{ref-sederberg2008context}
Sederberg, P. B., Howard, M. W., \& Kahana, M. J. (2008). A
context-based theory of recency and contiguity in free recall.
\emph{Psychological Review}, \emph{115}(4), 893.

\bibitem[\citeproctext]{ref-talmi2019retrieved}
Talmi, D., Lohnas, L., \& Daw, N. D. (2019). A retrieved context model
of the emotional modulation of memory. \emph{Psychological Review},
\emph{126}(4), 455.

\bibitem[\citeproctext]{ref-yonelinas2019contextual}
Yonelinas, A. P., Ranganath, C., Ekstrom, A. D., \& Wiltgen, B. J.
(2019). A contextual binding theory of episodic memory: Systems
consolidation reconsidered. \emph{Nature Reviews Neuroscience},
\emph{20}(6), 364--375.

\bibitem[\citeproctext]{ref-zaromb2006temporal}
Zaromb, F. M., Howard, M. W., Dolan, E. D., Sirotin, Y. B., Tully, M.,
Wingfield, A., \& Kahana, M. J. (2006). Temporal associations and
prior-list intrusions in free recall. \emph{Journal of Experimental
Psychology: Learning, Memory, and Cognition}, \emph{32}(4), 792--804.
\url{https://doi.org/10.1037/0278-7393.32.4.792}

\end{CSLReferences}






\end{document}
